% Product Management — Full Midterm (all week-8 content, word-packed)
% Compile: pdflatex midterm-print-all.tex

\documentclass[9pt,a4paper]{article}
\usepackage[utf8]{inputenc}
\usepackage[T1]{fontenc}
\usepackage[margin=1.1cm,top=0.9cm,bottom=0.9cm]{geometry}
\usepackage{enumitem}
\setlist{nosep,leftmargin=*,itemsep=0.15em,topsep=0.2em}
\usepackage{booktabs}
\usepackage{parskip}
\setlength{\parskip}{0.15em}
\usepackage{xcolor}
\usepackage{microtype}
\usepackage{multicol}
\usepackage[colorlinks=true,linkcolor=eublue,urlcolor=eublue]{hyperref}
\definecolor{eublue}{RGB}{0,51,153}
\definecolor{charcoal}{RGB}{51,51,51}
\definecolor{softgrey}{RGB}{248,248,252}

\begin{document}
\raggedright
\begin{center}
{\large\bfseries\textcolor{eublue}{Product Management — Midterm (All-in-One)}}\\[0.1em]
\footnotesize Mar 16 \quad|\quad Open notes
\end{center}
\vspace{-0.2em}

\noindent\fcolorbox{eublue}{white}{\parbox{\dimexpr\textwidth-2\fboxsep-2\fboxrule}{\footnotesize\textbf{How to use:} P1 = Exam day + Cheat sheet. P2+ = Model answers (Yiami, Dacia, Moley) + Extra practice (Dacia/Yiami/Moley style + mixed). Exam: read case first, underline clues, four questions; template every answer.}}
\vspace{0.25em}

\fcolorbox{eublue}{softgrey}{\parbox{\dimexpr\textwidth-2\fboxsep-2\fboxrule}{\footnotesize
\textbf{\color{eublue}EXAM DAY.} \textbf{Template:} (1) Name concept (2) Define one sentence (3) Three evidence from case (4) Conclude. \textbf{Four questions:} Cost/diff/response? Mass/custom? Create/capture value? Explore/exploit? Read case first; underline clues; ask four questions. \textbf{Time:} 5 min read; then template per question; last 15 min check concept+definition+evidence+conclusion.
}}
\vspace{0.25em}

% ========== CHEAT SHEET IN TWO COLUMNS ==========
\begin{multicols}{2}
\footnotesize

\textbf{\color{eublue}1. Make/Stock/Order.} STOCK = before orders. Make to Stock = before (bread, McD pre-made). Make to Order = after (piano, Made for You). Assemble to Order = parts before, assemble after (Dell, BMW). Memory: Stock = there; Order = you asked.

\textbf{\color{eublue}2. Lean vs.\ Design.} Design = observe, empathize, find problem. Lean = test \textbf{pay} before building. Stages: Empathize$\to$Define$\to$Ideate$\to$Prototype$\to$Test.

\textbf{\color{eublue}3. Innovation.} Sustaining/Disruptive: similar exist? YES$\to$Sust. NO or cheaper/new market$\to$Disrupt. Radical/Incremental: tech existed? YES$\to$Incr. NO$\to$Rad. Yiami/Dacia=Sust.+Incr.; Moley=Disrupt.+Rad.; Netflix=Disrupt.+Incr.; iPhone=Disrupt.+Rad.

\textbf{\color{eublue}4. Extreme users vs.\ Personas.} Extreme = use intensely; reveal problems (IDEO cart). Personas = who you design for; define after observing.

\textbf{\color{eublue}5. S-Curve vs.\ Innovator's Dilemma.} S-Curve = tech over time (slow$\to$rapid$\to$plateau). Innovator's Dilemma = why firms fail (Kodak, Blockbuster). Graph vs.\ trap.

\textbf{\color{eublue}6. Moore's Chasm.} Innovators = broken, tinker. Early Adopters = works, vision. Early Majority = proven. Late/Laggards = last. Memory: broken don't care / works see vision.

\textbf{\color{eublue}7. Value create/capture.} Create = useful. Capture = money. Three: Brand (Chanel); Market access (Takeda); IP/Patents (Samsung). Need both.

\textbf{\color{eublue}8. Break even.} Fixed$\div$(Price$-$Variable). Example: 6000$\div$(40$-$20)=300. Bracket first.

\textbf{\color{eublue}9. Learning effect.} More volume $\to$ better/cheaper. Dacia: (1) CMF-A shared $\to$ volume (2) Low price $\to$ demand $\to$ volume.

\textbf{\color{eublue}10. Economies of scale.} More units $\to$ lower cost/unit (fixed spread). Vs.\ learning = repetition.

\textbf{\color{eublue}11. Strategies.} Cost (Dacia, Ryanair). Differentiation (Apple, Chanel, Moley). Response (Zara, FedEx). Clues: cheap$\to$Cost; premium$\to$Diff.; fast$\to$Response.

\textbf{\color{eublue}12. BMC.} Roof: Value Prop. Right: Segments, Channels, Relationships. Left: Resources, Activities, Partners. Floor: Revenue, Cost. Specific names; one story.

\textbf{\color{eublue}13. Design — Yiami.} Order: Observe$\to$Personas$\to$Needs$\to$Ideate$\to$Prototype$\to$Test real users. Mistakes: personas first; no needs; straight to functionality; tested with devs. Good: ``last time...'' Bad: ``what features?''

\textbf{\color{eublue}14. Four questions.} Cost/diff/response? Mass/custom? Create/capture? Explore/exploit?

\textbf{\color{eublue}15. Manufacturing.} Craftsman: small, general, high skill, MTO, wide, low inv. Factory: large, specialized, low skill, MTS, narrow, high inv. Craft=diff.; Factory=cost.

\textbf{\color{eublue}16. House of Quality.} 3 requirements, weight 5/3/1. Rate components 5/3/1. Score=$\sum$(weight$\times$rating). Highest first.

\textbf{\color{eublue}17. Explore vs.\ Exploit.} Explore = new; risky. Exploit = what works; safe. Need both. Lean = test before invest.

\textbf{\color{eublue}18. Goods vs.\ services.} Heterogeneity; info asymmetry; customer involvement.

\textbf{\color{eublue}Key cases.} Moley: Disrupt+Rad, diff, innovators, R\&D. Dacia: Sust+Incr, cost, CMF-A, Zity, learning. Yiami: Sust+Incr, DT mistakes, B2B, innovators. Takeda, Xerox, IBM, McD, Blockbuster, Kodak, Zara, IDEO, Chanel.

\end{multicols}

\vspace{0.3em}
\footnotesize
\noindent\textbf{\color{eublue}More concepts.}\\
\textbf{Value Chain (Porter).} Primary: Inbound logistics, Operations, Outbound logistics, Mktg\&Sales, Service. Support: Infrastructure, HRM, Technology dev, Procurement. Dacia=Operations+Procurement; Apple=Tech+Mktg; Zara=Operations speed+Inbound.\\
\textbf{First mover.} Advantages: brand, loyalty, learning curve, network effects. Disadvantages: education cost, tech uncertainty, fast follower copies. Holds: high switching cost, strong network, learning. Fails: tech changes fast, no switching cost, follower copies. Moley: first robot kitchen = brand+patents; £50k copy = advantage can disappear.\\
\textbf{Network effects.} More users $\to$ more value. Direct: WhatsApp, Yiami (more users = more friends). Indirect: Airbnb (hosts/guests), Amazon (sellers/buyers). Once critical mass, hard to disrupt; crossing chasm = get to early majority then growth accelerates.\\
\textbf{Platform vs product.} Product: compete on features, direct revenue, you control (Dacia sells cars). Platform: others create value, revenue from ecosystem (commissions, ads), network effects (App Store 30\%). IBM open source = platform. Office=product; Azure=platform.\\
\textbf{Goods vs services (4).} Intangibility (can't touch before buy $\to$ brand matters). Heterogeneity (varies each time). Inseparability (produced+consumed at once $\to$ hard to scale). Perishability (can't store $\to$ yield mgmt).\\
\textbf{Bowling pin (cross chasm).} One niche first; dominate; use as reference for next. Mainstream wants proof. Yiami: Milan$\to$Paris$\to$Utrecht$\to$Munich. Moley: London enthusiasts$\to$NY$\to$cities$\to$mid-range.\\
\textbf{Ambidextrous.} Balance explore + exploit. Hard: explore=risk/freedom; exploit=efficiency. Do: separate divisions, different metrics, ring-fence budget. Microsoft: Office/Windows + Azure/OpenAI. Blockbuster: only exploit, died.\\
\textbf{Lean Canvas.} Problem, Solution, Unfair advantage, Key metrics (vs BMC's Partners/Activities). For early-stage; BMC for established. Same question: does the system make sense?\\
\textbf{CSDA.} Customer-focused, Strategic, Data-driven, Action-oriented. Role of PM = this.\\
\textbf{Quotes.} ``Creating $\neq$ capturing value'' (Teece). ``Less info $\to$ more emotion.'' ``Community before product.'' ``Build it and they will come = mistake.'' ``Testing = investing.'' ``Connection needs repetition'' (Yiami). ``People don't know what they want'' (Jobs). ``Great design can't save weak business model.''\\
\textbf{Exam traps.} (1) Moley sustaining? No — Thermomix blends, MRK cooks; disruptive. (2) Zara differentiation? No — response (speed). (3) MTO vs MTS: stock = before. (4) Innovators = broken; early adopters = works. (5) Canvas: specific names (Renault not ``partner''). (6) S-Curve = tech graph; Innovator's Dilemma = firm trap. (7) Name concept first. (8) Short structured beats long rambling.

\textbf{Willingness to pay.} Max customer will pay. No pay = no business; free $\neq$ customer; enthusiasm $\neq$ WTP; only test = ask for money. Test: pre-orders, landing page + buy button, manual offer, ``will you pay X now?'' Social app = built all, tested nothing. Won't pay: (1) problem not painful (2) solution not good enough (3) price too high. Yiami might pay: unis (retention budget), new city (acute pain), companies (onboarding).\\
\textbf{Blue ocean.} New uncontested space vs red = existing fight. Create new demand. Cirque (circus+theatre), Wii (non-gamers), Moley (autonomous cooking). Exam: Moley = blue; Dacia = red (cheap cars).\\
\textbf{Patents.} Prevent copy + prevent others blocking you. Portfolio: non-core, leverage, thickets. Solar case: A vs B — license / design around / acquire / negotiate. Moley: new tech, no patent = copy; recipe IP. Open innovation: IBM gave software for hardware; Moley should NOT open — tech IS product.\\
\textbf{Ansoff.} Penetration = existing product+market (Dacia more Springs). Market dev = existing product new market (Dacia new countries). Product dev = new product existing market (Dacia Cargo). Diversification = new+new (highest risk). Penetration/dev = exploit; diversification = explore.\\
\textbf{Pricing.} Cost-plus = cost+margin (Dacia). Value-based = what customer thinks worth (Apple). Penetration = cheap first, volume (Dacia). Skimming = expensive first then lower (Moley £248k, iPhone). Dynamic = demand-based (airlines, Uber).\\
\textbf{Stakeholders.} Internal: eng (excellence), sales (close deals), mktg (story), finance (profit), CEO (growth). External: customers, partners, investors, regulators. PM = balance all; customer focus. Conflict: sales wants feature, eng says 3mo — PM decides.\\
\textbf{Jobs to be done.} Customers hire product to do job. Drill $\to$ hole $\to$ picture $\to$ beautiful home. Yiami: job = connection not matching. Moley: job = effortless good food (recipe library, quality). Five whys $\to$ real job.\\
\textbf{Product life cycle.} Intro: low sales, R\&D, innovators (Moley). Growth: rising, mktg (Tesla). Maturity: plateau, efficiency (Dacia). Decline: harvest/exit. Budget by stage: intro = R\&D; growth = mktg; maturity = efficiency.\\
\textbf{Supply chain.} Upstream/downstream. Vertical integration = control, high fixed (Apple chips). Outsourcing = flex, low fixed (Dacia–Renault). JIT = materials when needed (Toyota). Just in case = buffer (post-COVID).\\
\textbf{How to use notes.} (1) Write four magic questions at top. (2) Read case once, no writing. (3) Read again, underline clues (cheap$\to$cost; premium$\to$diff; fast$\to$response; existing$\to$sust; new$\to$disrupt; existing tech$\to$incr; new tech$\to$rad). (4) Find relevant section before each answer. (5) Template every time. (6) Specific names. (7) Use extra time; check all. (8) If blank: four questions $\to$ answer about company $\to$ build from there.

\vspace{0.2em}
\newpage
% ========== MODEL ANSWERS (DENSE) ==========
\footnotesize
\section*{\color{eublue}Model Answers (Yiami, Dacia, Moley)}

\textbf{Yiami Q1.} Sustaining (Tinder/Meetup exist). Incremental (AI, geo, chat exist).\\
\textbf{Yiami Q2.} Personas$\to$Observing$\to$Functionality$\to$Tech test$\to$Launch. Mistakes: (1) Personas before observing (2) No needs What/How/Why (3) No ideation (4) Tested Reddit devs not target users. HoQ: Chatting 1, Compatible 3, Events offline 5. Scores: DB filters 41 first, Matching 31, Geo 31, Event 29, Search 21, ID 17, Conv 13.\\
\textbf{Yiami Q3.} Innovators; not finished; target tech enthusiasts. Niche: unis Milan/Paris/Utrecht/Munich = beachhead.\\
\textbf{Yiami Q4.} Training: market/competitors; personas; value prop; product/demos; objections.

\textbf{Dacia Q1.} Cost leadership. Evidence: CMF-A Renault/Dongfeng/Nissan; China specialised workers; minimal features; simple marketing.\\
\textbf{Dacia Q2.} Specialised. Repetitive = lower wages = cost leadership. Skilled = whole craft = higher = undermines.\\
\textbf{Dacia Q3.} (1) Share CMF-A $\to$ volume $\to$ learning. (2) Low price $\to$ demand $\to$ volume $\to$ learning.\\
\textbf{Dacia Q4.} EoS = more output $\to$ lower per-unit fixed cost. Dacia: low price (volume); shared platform (volume).\\
\textbf{Dacia Q5.} Heterogeneity (services vary). Info asymmetry (brand signal). Customer involvement (low productivity).\\
\textbf{Dacia Q6.} VP: affordable EV, Zity. Activities: mfg, marketing. Resources: CMF-A. Partners: Renault, Dongfeng, Nissan, Zity. Cost: low fixed/marketing/variable (China). Rel: non-recurrent showrooms, recurrent Zity. Segments: cheap EV, Zity, Cargo, car-share. Channels: showrooms, Zity. Revenue: sales, sharing.

\textbf{Moley Q1.} Differentiation. Evidence: £248k; autonomous 5k recipes; enthusiasts; Thermomix €1k can't cook. Sustainable: R\&D; patents; recipe ecosystem; cheaper later.\\
\textbf{Moley Q2.} VP: first autonomous chef, 5k recipes, restaurant at home. Segments: wealthy, enthusiasts, chefs, hotels. Channels: Moley, showrooms, trade. Rel: premium, updates. Revenue: £248k, subs, maintenance. Resources: arms, recipe DB, AI, patents. Activities: R\&D, capture, software, mfg. Partners: chefs, suppliers, luxury. Cost: high R\&D, value driven.\\
\textbf{Moley Q3.} Intro/innovators. R\&D 35\%, Tech design 25\%, Market research 15\%, Production 10\%, Marketing 10\%, Sales 5\%.\\
\textbf{Moley Q4.} Disruptive (new market; Thermomix blends, MRK cooks). Radical (arms = new tech). Budget: radical+disrupt early $\to$ heavy R\&D (60\%); sustaining+incr like Dacia $\to$ more marketing.

\vspace{0.4em}
\section*{\color{eublue}Extra practice (professor-style)}

\textbf{Dacia — Porter Value Chain.} Shared platform $\to$ \textbf{Operations} (mfg cost cut; fixed shared). Inbound logistics (bulk, suppliers). Dacia = world-class efficiency without world-class investment (complementary asset).\\
\textbf{Dacia — Low price, OK margins.} EoS (fixed spread over volume); learning effect (more cars $\to$ cheaper); low variable (China, old Renault tech). Cost leadership: costs so low that low price still profits.\\
\textbf{Dacia — Sales vs.\ Zity (goods vs.\ services).} (1) Heterogeneity: Spring = consistent; Zity = varies (dirty car, slow app). Different QC. (2) Customer involvement: sale = done; Zity = customer finds/unlocks/drives/parks — design for that friction.

\textbf{Yiami — Multiple needs (deep friends vs.\ network vs.\ hobbies).} Multiple segments; one product can't serve all. Use House of Quality: weight requirement that drives download; prioritise (e.g.\ offline events = core for real connection).\\
\textbf{Yiami — ``Infinite strangers made me lonelier.''} Paradox of choice; design: limit connections (scarcity); quality match over quantity; re-matching = depth. Challenge = intimate/selective vs.\ infinite scroll.\\
\textbf{Yiami — B2B university pivot.} (1) WTP: unis pay (retention budget); users don't. (2) Chicken-egg: unis = concentrated new students, same problem. (3) Teece: unis have distribution, trust, motivation; Yiami captures via their complementary assets.

\textbf{Moley — Can't peel/dice.} Innovators stage (tolerate limits). Blocks chasm (mainstream won't pay £248k for half-job). Do: R\&D on full autonomy; focus on enthusiasts/chefs meanwhile.\\
\textbf{Moley vs.\ Thermomix — Same market?} No. Thermomix = assist cook (weigh, blend); MRK = replace cook (autonomous). Different value prop. MRK = disruptive, new category; must educate market; heavy R\&D not mass marketing.\\
\textbf{Moley — Competitor at £50k.} Disruptive from below. Innovator's Dilemma: current customers want more; instinct = premium; £50k gets different segment. S-curve: cheap can improve fast. Respond: patents; recipe ecosystem (lock-in); consider mid-range model (founder hinted).

\textbf{Learning vs.\ EoS.} EoS = fixed costs spread over more units (math). Learning = repetition $\to$ better/cheaper (experience). Example EoS: Dacia CMF-A shared $\to$ fixed over millions. Example learning: Dacia low price $\to$ volume $\to$ experience.\\
\textbf{``Market wasn't ready.''} Professor right: validate before building (lean). ``Not ready'' often = didn't test WTP/timing. Social app = travelled hoping; lean = find no WTP week one. Test timing; take responsibility.\\
\textbf{Great design, bad business model.} Agree. Folding bike = good design, bad fit, failed. Xerox PARC = mouse/GUI/Ethernet, no commercialisation; Apple/Microsoft captured. Design = one BMC box; need all nine. Design + business model both needed.

\vspace{0.3em}
\noindent\rule{0.4\textwidth}{0.3pt}\quad \small\textcolor{charcoal}{JZ — UCM Product Management.}
\end{document}
